\centerline{\bf \Huge ДЗ. Перебор случаев+Обратный ход}
\centerline{\bf \Huge }

\begin{flushright}
        \scriptsize{\textbf{Ответы без объяснения решения оцениваются в 0 баллов!}}
\end{flushright}

Каждая задача стоит \textbf{3 балла}, но при этом 1 и 2 задачи \textbf{обязательные}! 

За неправильное выполнение 1 и 2 задачи можно получить \textbf{отрицательное} количество баллов!

\problem{1!} Герман загадал число, затем прибавил к нему 5 и после разделил на 3. Какое число загадал Герман, если после его действий у него получилось число 7?

\problem{2!} Робот передвигается вдоль прямой вперёд на 110м или назад на 70м. Может ли он удалиться от начальной точки ровно на 777м?

\problem{3} Назовём дату «идеальной», если сумма первых четырёх цифр равна сумме последних четырёх. Пример «идеальной» даты 07.10.2024. Перечислите все «идеальные» даты в 2024 году.

\problem{4}  На прямой отметили несколько точек. После этого между каждыми двумя соседними точками отметили ещё по две точки. Такую операцию повторили ещё дважды (всего 3 раза). В результате на прямой оказалось отмечено 244 точки. Сколько точек было отмечено первоначально?

